\section{Phạm vi so sánh và nền tảng liên quan}

N-Queens là bài toán khả thi: với một $N$ cố định, đầu ra cần là một bàn cờ
hợp lệ hoặc chứng nhận không có nghiệm. Cách đặt vấn đề này quan trọng khi
so sánh solver. SAT và CP-SAT có thể trả về phép gán; MIP cũng chỉ cần một
nghiệm khả thi, không có giá trị mục tiêu để tối ưu. Do đó thời gian trong
báo cáo chỉ nói về việc tìm một nghiệm hoặc kết thúc với kết luận UNSAT trong
giới hạn đã cấu hình, không nói về số nghiệm hay chất lượng tối ưu.

Ba họ mô hình có cùng tập nghiệm nhưng trình bày thông tin cho solver theo
những cách khác nhau. SAT chuyển các nhóm đếm thành CNF; PySAT cung cấp
giao diện tạo ràng buộc cardinality và gọi Glucose3~\cite{pysat-paper,pysat-docs}.
CP giữ biến nguyên theo hàng và ba nhóm \texttt{AllDifferent}, đúng với mô
hình N-Queens của OR-Tools~\cite{ortools-nqueens}. MIP giữ biến nhị phân theo
ô cùng các tổng tuyến tính. Vì vậy số biến của SAT và MIP không so trực tiếp
với số biến CP: biến nguyên CP có miền $N$ giá trị, còn các biến Boolean SAT
có thể bao gồm biến phụ của encoding.

Đối với SAT, sự khác biệt giữa các mã hóa AMO có hai tầng. Tầng thứ nhất là
chi phí dựng công thức: số biến phụ, số clause và độ dài clause. Tầng thứ
hai là hiệu quả tìm kiếm sau khi công thức đã được đưa vào Glucose3: mức độ
lan truyền, học clause và lựa chọn biến. Sequential Counter giảm kích thước
clause theo nhóm so với Pairwise~\cite{sinz2005}, nhưng điều đó không tự nó
bảo đảm thời gian giải thấp hơn. Báo cáo vì thế trình bày cả cấu trúc CNF,
thời gian và tỷ lệ hoàn thành; một con số thời gian đơn lẻ không đủ.

Phạm vi thực nghiệm là các $N$ xuất hiện trong bộ dữ liệu đã công bố:
$1,2,3,4,8,16,32,64,128,256,512,1000$. Không phải mọi phương pháp đều
được chạy ở mọi $N$. Những ô không có dữ liệu được giữ là không chạy,
không diễn giải là thất bại. Các so sánh có ý nghĩa nhất nằm trong cùng
kích thước và cùng giới hạn thời gian; riêng Gurobi ở $N=64,128,256$ được
bổ sung ở phiên Colab khác nên cần đọc với lưu ý về môi trường.
